\ChapterImageStar[cap:caracterizacion]{Caracterización de las metodologías y tecnologías}{./images/fondo.png}\label{cap:caracterizacion}
\mbox{}\\

En esta sección se presentan los principales referentes metodológicos y tecnológicos identificados durante el proceso de revisión 
sistemática de la literatura. Su propósito es analizar las metodologías, arquitecturas y herramientas de seguridad más relevantes 
encontradas en los estudios seleccionados, con el fin de comprender sus características, capacidades y posibles aplicaciones dentro del 
contexto de la infraestructura de computación en la nube del grupo GRID.\@

\section{Caracterización de Referentes Metodológicos}

Esta etapa se centra principalmente en la caracterización de los referentes metodológicos identificados durante el proceso de SMS.\@ Su 
objetivo es analizar los atributos, capacidades y alcances de cada metodología encontrada, proporcionando un análisis detallado que 
permita comprender su aplicabilidad en el contexto de la seguridad de la infraestructura en la nube del grupo GRID.\@ Esta caracterización 
constituye un insumo fundamental para la posterior selección de metodologías en el desarrollo del presente trabajo de grado.

\subsection{Estándar ISO/IEC 27001}

\subsubsection{Contexto y relevancia del estándar ISO/IEC 27001}

De acuerdo con \citeauthor{rajesh2024unified}, la norma ISO/IEC 27001 constituye un marco de cumplimiento regulatorio fundamental en 
entornos de computación en la nube, especialmente en un contexto donde la adopción de servicios cloud continúa creciendo y las 
preocupaciones relacionadas con la seguridad de la información representan uno de los principales desafíos para las organizaciones. 
Su relevancia radica en proporcionar lineamientos orientados a garantizar la gestión de la seguridad de la información bajo un modelo de 
responsabilidad compartida entre usuarios, proveedores y demás actores involucrados en la infraestructura cloud.

Asimismo, el estudio desarrollado por \citeauthor{tissir2021cybersecurity} resalta que este estándar permite definir necesidades de 
gestión y control que aún presentan limitaciones dentro de la literatura actual, particularmente en entornos dinámicos y distribuidos 
como los de computación en la nube.

\subsubsection{Alcance prestacional}

Con base en los estudios analizados, los principales aspectos funcionales de la ISO/IEC 27001 se centran en los siguientes elementos:

\begin{itemize}
	\item \textbf{Gestión integral de riesgos:} Proporciona un marco para identificar, evaluar y mitigar riesgos relacionados con la seguridad de la información, considerando la interacción entre personas, procesos y tecnologías~\cite{tissir2021cybersecurity}.
	
	\item \textbf{Estandarización y auditoría:} Funciona como un referente internacional para la realización de auditorías y procesos de verificación del cumplimiento de controles de seguridad en infraestructuras tecnológicas~\cite{gambo2025navigating}.
	
	\item \textbf{Marco de políticas y controles:} Establece lineamientos orientados a la definición de políticas, procedimientos y controles de seguridad que permitan fortalecer la protección de los activos de información y mejorar la madurez organizacional~\cite{tissir2021cybersecurity}.
\end{itemize}

\subsubsection{Aplicabilidad en el contexto de infraestructura en la nube}

La aplicación de la norma ISO/IEC 27001 en entornos de computación en la nube permite abordar problemáticas relacionadas con la auditoría, 
el aseguramiento de servicios y la gestión de riesgos en infraestructuras cloud. Diversos estudios destacan que este estándar facilita la 
transparencia entre los actores involucrados y contribuye a medir el nivel de madurez de las organizaciones frente a la gestión de la 
seguridad de la información~\cite{rajesh2024unified,tissir2021cybersecurity}.

De igual manera, la literatura señala que su implementación debe considerar factores como la disponibilidad de personal especializado, 
los recursos económicos y las capacidades tecnológicas de la organización, especialmente en contextos institucionales y 
académicos~\cite{gambo2025navigating}.

\subsubsection{Sustento científico vinculado}

El sustento científico asociado a la ISO/IEC 27001 dentro de esta investigación se encuentra respaldado por los SPS 084, 154 y 025, los 
cuales presentan enfoques relacionados con auditoría de seguridad, gestión de riesgos y cumplimiento regulatorio en entornos cloud. En 
conjunto, estos estudios posicionan a la ISO/IEC 27001 como un componente fundamental para fortalecer la resiliencia operativa, la 
protección de activos de información y la gestión de la ciberseguridad en infraestructuras de computación en la 
nube~\cite{rajesh2024unified,tissir2021cybersecurity}.

%####################################################################################################################################

\subsection{NIST Cybersecurity Framework (CSF)}

\subsubsection{Contexto y relevancia del \NISTCSF}

El marco de ciberseguridad del \textit{National Institute of Standards and Technology (\NISTCSF)} constituye un referente internacional 
orientado a la definición de lineamientos y mejores prácticas para la gestión de la ciberseguridad dentro de las organizaciones. De 
acuerdo con~\cite{tissir2021cybersecurity}, este marco proporciona un lenguaje común que facilita la maduración de procesos de seguridad y 
fortalece la gestión de riesgos en entornos tecnológicos distribuidos.

Asimismo,~\cite{rajesh2024unified} caracterizan al \NISTCSF como un estándar técnico esencial para establecer procedimientos de auditoría y 
verificación de controles de seguridad en entornos cloud. Su relevancia también se fundamenta en su capacidad para responder a las 
necesidades de gestión derivadas de paradigmas emergentes como la computación distribuida y la computación en la nube.

Por otra parte, el estudio desarrollado por~\cite{ulloa2025integral} destaca que este marco no solo resulta aplicable en servidores e 
infraestructuras tradicionales, sino también en ecosistemas de dispositivos interconectados, donde la protección continua de los activos y 
canales de comunicación representa un aspecto crítico para prevenir vulnerabilidades y ataques.

\subsubsection{Alcance prestacional}

A partir de la síntesis de los estudios analizados, los principales aspectos funcionales del \NISTCSF se centran en los siguientes 
elementos:

\begin{itemize}
	\item \textbf{Seguridad desde el diseño e integridad:} El marco promueve la protección frente a accesos no autorizados y modificaciones maliciosas de los datos, garantizando la integridad de los activos dentro de entornos interconectados~\cite{ulloa2025integral}.
	
	\item \textbf{Adaptabilidad y maduración organizacional:} Diversos estudios destacan que el \NISTCSF permite establecer un lenguaje común entre las organizaciones, facilitando la maduración de procesos de seguridad y la gestión de riesgos~\cite{tissir2021cybersecurity,rajesh2024unified}.
	
	\item \textbf{Estandarización y apoyo al cumplimiento:} El marco funciona como un referente internacional que puede complementarse con estándares como ISO/IEC 27001 para fortalecer procesos de auditoría, cumplimiento y gobernanza de la seguridad de la información~\cite{gambo2025navigating}.
\end{itemize}

\subsubsection{Aplicabilidad en el contexto de infraestructura en la nube}

La aplicabilidad del \NISTCSF dentro del contexto de infraestructura en la nube se orienta principalmente a tres aspectos identificados en 
la literatura:

\begin{itemize}
	\item \textbf{Auditoría y confianza:} El \NISTCSF facilita la implementación de procesos continuos de auditoría y verificación de controles, fortaleciendo el cumplimiento técnico y operativo en entornos cloud~\cite{rajesh2024unified}.
	
	\item \textbf{Protección de dispositivos e infraestructura:} El marco contribuye a asegurar los canales de comunicación entre dispositivos y servicios en la nube, favoreciendo la protección de los flujos de información y la integridad de los activos tecnológicos~\cite{ulloa2025integral}.
	
	\item \textbf{Evaluación del nivel de madurez:} El \NISTCSF permite evaluar la capacidad de respuesta y madurez organizacional frente a amenazas emergentes, apoyando la gestión integral de riesgos y vulnerabilidades~\cite{tissir2021cybersecurity,gambo2025navigating}.
\end{itemize}

\subsubsection{Sustento científico vinculado}

El sustento científico asociado al \NISTCSF dentro de esta investigación se encuentra respaldado por los SPS 025, 084, 134 y 154. Los 
estudios analizados posicionan este marco como un enfoque orientado a fortalecer la gobernanza de la seguridad, la estandarización de 
criterios y la maduración de procesos organizacionales en contextos tecnológicos heterogéneos~\cite{rajesh2024unified,ulloa2025integral}.

%####################################################################################################################################

\subsection{ISO/IEC 27017}

\subsubsection{Contexto y relevancia del estándar ISO/IEC 27017}

El estándar ISO/IEC 27017 presentó una presencia limitada dentro del proceso de SMS desarrollado en esta investigación, siendo 
identificado principalmente en el estudio realizado por~\cite{tissir2021cybersecurity}. Dicho trabajo posiciona este estándar como un 
componente relevante para abordar los desafíos asociados al paradigma emergente de la computación en la nube, especialmente en aspectos 
relacionados con la gestión de la ciberseguridad y la responsabilidad compartida entre proveedores y consumidores de servicios cloud.

La relevancia de la ISO/IEC 27017 radica en que proporciona lineamientos y controles específicos orientados a entornos de computación en la 
nube, permitiendo complementar otros marcos y estándares de seguridad de la información para fortalecer la gestión de riesgos, amenazas y 
vulnerabilidades dentro de infraestructuras cloud~\cite{tissir2021cybersecurity}.

\subsubsection{Aplicabilidad en el contexto de la infraestructura en la nube}

El framework conceptual propuesto por~\cite{tissir2021cybersecurity}, aunque no se enfoca específicamente en entornos universitarios, 
resalta que la ISO/IEC 27017 contribuye al fortalecimiento de la madurez organizacional y a la protección de los activos de información 
presentes en infraestructuras de computación en la nube.

Asimismo, el estudio destaca que cualquier infraestructura cloud se encuentra expuesta a amenazas, vulnerabilidades y ataques, razón por 
la cual la integración de este estándar con otros marcos de referencia y normativas permite establecer enfoques de seguridad más robustos 
y alineados con las necesidades de protección de los entornos distribuidos.

\subsubsection{Sustento científico vinculado}

El sustento científico asociado a la ISO/IEC 27017 dentro de esta investigación se encuentra respaldado principalmente por el SPS 084 
correspondiente al estudio desarrollado por~\cite{tissir2021cybersecurity}. Los autores proponen un modelo conceptual basado en análisis 
semántico que integra estándares como ISO/IEC 27017, ISO/IEC 27001 y el \NISTCSF, enfatizando la importancia de la madurez organizacional, 
la gestión de riesgos y la implementación de controles de seguridad orientados a la prevención de amenazas y vulnerabilidades en entornos cloud.

%####################################################################################################################################

\subsection{Zero-Trust Architecture (\ZTA)}

\subsubsection{Contexto y relevancia de la arquitectura Zero-Trust Architecture (\ZTA)}

La arquitectura de ciberseguridad \textit{Zero-Trust Architecture (\ZTA)} constituye un modelo orientado a fortalecer la protección de los 
activos de información mediante el principio de ``nunca confiar, siempre verificar''. De acuerdo con~\cite{ali2025zta}, este enfoque trata 
la identidad como el nuevo perímetro de seguridad, realizando procesos continuos de autenticación y validación tanto de usuarios como de 
dispositivos antes de conceder acceso a cualquier recurso dentro de la infraestructura.

Asimismo,~\cite{stojalowski2024zero} describen el \ZTA como un modelo de seguridad diseñado para proteger datos y servicios sin importar su 
ubicación, incluyendo la prevención de movimientos laterales dentro de la infraestructura. Los autores destacan que este enfoque se 
fundamenta en pilares relacionados con la seguridad de usuarios, dispositivos, aplicaciones, redes, datos, visibilidad, analítica y 
automatización, permitiendo fortalecer el control sobre los activos procesados en entornos distribuidos y cloud.

\subsubsection{Alcance prestacional}

Con base en los estudios analizados sobre Zero-Trust Architecture (\ZTA), se identifica que este modelo de ciberseguridad protege los 
activos de información mediante mecanismos continuos de autenticación, validación y monitoreo de usuarios y dispositivos. De acuerdo 
con~\cite{bennaceur2026zero}, el \ZTA permite proteger información crítica frente a comportamientos maliciosos, incorporando esquemas 
orientados a detectar y restringir actividades sospechosas dentro de la infraestructura.

Asimismo,~\cite{stojalowski2024zero} destacan que el concepto de Zero-Trust se encuentra respaldado por la publicación especial NIST0
SP 800-207, consolidándose como un enfoque orientado a la protección de datos, servicios y activos de información empresariales. Los 
autores señalan que su alcance se fundamenta en la implementación de controles distribuidos sobre distintos componentes de la 
infraestructura, incluyendo usuarios, dispositivos, aplicaciones, cargas de trabajo, redes, datos, visibilidad, analítica y automatización.

De igual manera, la literatura resalta que el principio de ``nunca confiar, siempre verificar'' se aplica mediante procesos constantes de 
autenticación, inspección de tráfico, generación de logs y monitoreo de actividades, independientemente de si el acceso proviene desde 
dentro o fuera de la infraestructura organizacional. Esto permite fortalecer la administración y protección de entornos de computación en 
la nube de manera más segura y controlada~\cite{stojalowski2024zero}.

\subsubsection{Aplicabilidad en la infraestructura en la nube}

La implementación de la arquitectura Zero-Trust en infraestructuras de computación en la nube requiere un proceso gradual de transición 
desde modelos tradicionales basados en perímetros hacia enfoques centrados en identidad y control continuo~\cite{stojalowski2024zero}. 
En este contexto, la literatura identifica principios fundamentales para su adopción:

\begin{itemize}
	\item Garantizar el acceso seguro a los recursos independientemente de su ubicación.
	
	\item Aplicar estrategias de mínimos privilegios y controles estrictos de acceso.
	
	\item Inspeccionar y monitorear continuamente el tráfico y los registros generados dentro de la infraestructura.
\end{itemize}

Estos principios permiten fortalecer el control sobre la superficie de protección y asegurar que los usuarios únicamente accedan a los 
recursos estrictamente necesarios para el desarrollo de sus actividades~\cite{stojalowski2024zero}.

\subsubsection{Sustento científico vinculado}

El sustento científico asociado a la arquitectura Zero-Trust dentro de esta investigación se encuentra respaldado por los SPS 151, 175 y 
180. Los estudios analizados presentan simulaciones y casos prácticos que evidencian la capacidad del modelo \ZTA para fortalecer la 
seguridad, estabilidad y resiliencia de infraestructuras tecnológicas frente a amenazas emergentes en entornos públicos y 
privados~\cite{bennaceur2026zero}.

%####################################################################################################################################

\section{Caracterización de Referentes Tecnológicos}

Posterior al establecimiento de los referentes metodológicos, esta fase se enfoca en la caracterización de los referentes tecnológicos. A 
partir de los hallazgos identificados en la literatura, se analizan las tecnologías desde sus capacidades de detección, mitigación y 
compatibilidad, permitiendo validar la viabilidad operativa de sus funciones frente a los requerimientos de seguridad identificados en el 
entorno de infraestructura en la nube del grupo \GRID.\@

\subsection{Intrusion Detection and Prevention System (\IDPS)}

\subsubsection{Contexto y relevancia del IDPS}

Los sistemas Intrusion Detection and Prevention System (\IDPS) desempeñan un rol fundamental en la protección de redes e infraestructuras 
tecnológicas mediante la detección, análisis y prevención de actividades maliciosas. De acuerdo con~\cite{nasim2025idps}, estas tecnologías 
integran las capacidades de detección de un \IDS con las funciones de prevención de un \IPS, permitiendo identificar amenazas y ejecutar 
acciones orientadas a mitigar ataques dentro de la infraestructura.

Asimismo, los autores destacan que, aunque los \IDPS fueron desarrollados inicialmente para entornos de red tradicionales, actualmente 
enfrentan desafíos asociados a la naturaleza dinámica y cambiante de los entornos de computación en la nube~\cite{nasim2025idps}.

Por otra parte,~\cite{vaneeta2025idps} atribuyen a los \IDPS la capacidad de mejorar su desempeño mediante el uso de técnicas híbridas 
como \textit{signature-based}, \textit{anomaly-based} y \textit{hybrid-based}, las cuales permiten optimizar la detección de amenazas y 
adaptar los mecanismos de seguridad a infraestructuras ya existentes. En conjunto, la literatura resalta la versatilidad de estas 
tecnologías y su capacidad para responder tanto a amenazas conocidas como a comportamientos anómalos dentro de entornos cloud.

\subsubsection{Capacidad operativa}

Con base en la literatura analizada, las principales capacidades operativas de las tecnologías IDPS se clasifican en los siguientes 
enfoques:

\begin{itemize}
	\item \textbf{Signature-based Technique:} Este enfoque permite detectar amenazas mediante la comparación del tráfico de red con firmas previamente conocidas. Su principal ventaja radica en la alta precisión frente a ataques reconocidos y en la reducción de falsos positivos. Sin embargo, presenta limitaciones frente a amenazas nuevas o desconocidas que no han sido registradas previamente~\cite{nasim2025idps}.
	
	\item \textbf{Anomaly-based Technique:} Este método se basa en el análisis del comportamiento normal de usuarios y sistemas para identificar desviaciones o actividades sospechosas. Aunque permite detectar amenazas desconocidas, puede generar mayores tasas de falsos positivos debido a variaciones frecuentes en los patrones de comportamiento~\cite{nasim2025idps}.
	
	\item \textbf{Hybrid-based Technique:} Este enfoque combina técnicas basadas en firmas y detección de anomalías, permitiendo incrementar la efectividad del IDPS en entornos de computación en la nube. Su integración facilita la detección tanto de amenazas conocidas como de ataques emergentes y vulnerabilidades desconocidas~\cite{nasim2025idps}.
\end{itemize}

\subsubsection{Aplicabilidad en la nube}

Según~\cite{adiwal2023cloud}, los proveedores de servicios cloud requieren la implementación de tecnologías \IDPS como parte fundamental 
de la protección de sus infraestructuras, siempre que estas se configuren de acuerdo con las necesidades y características específicas del 
entorno.

No obstante,~\cite{nasim2025idps} señalan que estas tecnologías pueden presentar limitaciones relacionadas con el consumo de recursos 
computacionales y la necesidad de realizar pruebas constantes sobre tráfico real y escenarios de amenazas dinámicas. A pesar de ello, los 
estudios coinciden en que una configuración y adaptación adecuada de los IDPS permite fortalecer significativamente la seguridad de 
infraestructuras de computación en la nube.

En consecuencia, la literatura posiciona a los \IDPS como tecnologías relevantes para la identificación, clasificación y mitigación de 
amenazas dentro de entornos cloud.

\subsubsection{Sustento científico vinculado}

El sustento científico asociado a las tecnologías \IDPS dentro de esta investigación se encuentra respaldado por los SPS 056, 129 y 137. 
Los estudios analizados evidencian que estas tecnologías poseen una amplia adopción en entornos de producción debido a su efectividad, 
versatilidad y disponibilidad de herramientas \textit{open-source} respaldadas por comunidades activas de desarrollo y 
soporte.

%####################################################################################################################################

\subsection{Security Information and Event Management (\SIEM)}

\subsubsection{Contexto y relevancia del Security Information and Event Management (\SIEM)}

Los sistemas \textit{Security Information and Event Management (\SIEM)} constituyen soluciones orientadas a la correlación, análisis y monitoreo de 
eventos de seguridad provenientes de múltiples fuentes dentro de una infraestructura tecnológica. De acuerdo con~\cite{coppolino2022siem}, 
estas herramientas permiten integrar información procedente de dispositivos de red, sistemas operativos, mecanismos de administración de 
identidades y plataformas de control de acceso, facilitando la identificación y análisis de incidentes de seguridad.

Asimismo,~\cite{hadi2024opensource} describen al SIEM como una tecnología especializada en la recolección, gestión y análisis de logs, 
desempeñando un rol crítico en la detección y administración de eventos de seguridad dentro de infraestructuras organizacionales. Los 
autores destacan que estas soluciones permiten centralizar grandes volúmenes de información para fortalecer los procesos de monitoreo y 
respuesta ante incidentes.

No obstante, la literatura también señala que las herramientas \SIEM \textit{open-source} pueden presentar desafíos relacionados con su 
complejidad de implementación y la necesidad de configuraciones avanzadas para adaptarse a los requerimientos específicos de cada 
organización. A pesar de ello, su flexibilidad, escalabilidad y bajo costo las convierten en alternativas atractivas para entornos 
institucionales y empresariales~\cite{hadi2024opensource}.

\subsubsection{Capacidad operativa}

Con base en la literatura analizada, la capacidad operativa de las tecnologías \SIEM se fundamenta principalmente en la monitorización, 
correlación y análisis de eventos de seguridad mediante el procesamiento de logs provenientes de diferentes componentes de la 
infraestructura~\cite{coppolino2022siem}.

Los estudios revisados destacan que los logs representan el elemento central dentro de este tipo de soluciones, ya que permiten 
identificar accesos autorizados y no autorizados, actividades sospechosas, movimientos de información y fallos dentro de los sistemas. De 
igual manera, el \SIEM facilita la priorización y análisis de incidentes por parte del personal encargado de la seguridad, fortaleciendo la 
capacidad de respuesta frente a amenazas y eventos críticos~\cite{coppolino2022siem}.

\subsubsection{Aplicabilidad en la nube}

De acuerdo con~\cite{coppolino2022siem}, las soluciones \SIEM se posicionan como herramientas estratégicas para el monitoreo y análisis de 
eventos dentro de infraestructuras cloud, proporcionando información oportuna para la detección y mitigación de ciberataques.

Asimismo, la literatura destaca que la adopción progresiva de servicios en la nube incrementa la necesidad de implementar mecanismos 
centralizados de monitoreo y correlación de eventos que permitan proteger la integridad y disponibilidad de los servicios críticos. En este 
contexto,~\cite{hadi2024opensource} señalan que los \SIEM constituyen componentes fundamentales dentro de estrategias de defensa multicapa, 
gracias a sus capacidades de análisis, correlación de eventos y administración de logs en entornos distribuidos.

Por consiguiente, el \SIEM se establece como una tecnología relevante para fortalecer los procesos de vigilancia, detección y gestión de 
incidentes dentro de infraestructuras de computación en la nube.

\subsubsection{Sustento científico vinculado}

El sustento científico asociado a las tecnologías \SIEM dentro de esta investigación se encuentra respaldado por los SPS 083 y 168. Los 
estudios analizados evidencian una amplia recomendación hacia la implementación de herramientas \SIEM en infraestructuras cloud, debido a 
su capacidad para monitorizar eventos, analizar actividades sospechosas y fortalecer la protección de los activos de información mediante 
mecanismos centralizados de gestión y correlación de logs.
%####################################################################################################################################
\subsection{Conclusiones de la caracterización}

Este proceso permitió sustentar la investigación sobre bases científicas y metodológicas, evitando un enfoque empírico y proporcionando 
argumentos sólidos y defendibles para la construcción de la propuesta.

Asimismo, la cantidad de estudios procesados permitió dimensionar la amplitud del ecosistema de la ciberseguridad en la nube, evidenciando 
múltiples enfoques, estrategias y tecnologías aplicables a infraestructuras cloud. Esto facilitó un proceso de selección y adaptación 
orientado a identificar los modelos y herramientas más adecuados para las necesidades del entorno del grupo GRID, favoreciendo el 
planteamiento de una solución técnica contextualizada y sustentada.

Para una descripción más detallada sobre la caracterización de los referentes metodológicos y tecnológicos analizados en esta etapa, se 
recomienda consultar el Apéndice~\ref{sec:caracterizacion-referentes}, donde se presenta la información complementaria asociada al proceso de análisis y selección 
realizado.